%% kapa: exam=1 round=35 year=2024 period=2 form=A subject=accounting subject_ko=회계학 seq=32 number=72
%% source: Q-Net 감정평가사 기출문제 (cst00309) artlSeq=5213229 file=raw/1008/5213229/2교시 문제지 원본.pdf
%% conversion: kordoc 4.24.51 -> tools/gen1.py (auto); review=proofread-sonnet; flags=tab,diagonal-header
\begin{question}{72}
(주)감평은 두 개의 제조부문 P1, P2와 두 개의 보조부문 S1, S2를 통해 제품을 생산하고 있다. S1과 S2의 부문원가는 각각 ￦60,000과 ￦30,000이다. 다음 각 부문간의 용역수수 관계를 이용하여 보조부문원가를 직접배분법으로 제조부문에 배분할 때 P2에 배분될 보조부문원가는? (단, S1은 기계시간, S2는 kW에 비례하여 배분한다.)
\par{\small\centering\begin{tabular}{|c|c|c|c|c|}\hline
제공 ＼ 사용 & \multicolumn{2}{c|}{제조부문} & \multicolumn{2}{c|}{보조부문} \tabularnewline\hline
 & P1 & P2 & S1 & S2 \tabularnewline\hline
S1 & 30기계시간 & 18기계시간 & 5기계시간 & 8기계시간 \tabularnewline\hline
S2 & 160kW & 240kW & 80kW & 50kW \tabularnewline\hline
\end{tabular}\par}
\begin{choices}
  \item ￦18,000
  \item ￦22,500
  \item ￦37,500
  \item ￦40,500
  \item ￦55,500
\end{choices}
\end{question}
\answer{4}
